\section{Purpose: reconstruct admissible magnetics from electrical samples}
An accurate interpolator can match two measured flux surfaces without making
them the gradient of any common magnetic potential. If those surfaces are
later differentiated, visually small inconsistencies may become large errors
in differential inductance. The research question here is therefore not only
how to minimize pointwise flux error, but how to reconstruct smooth saturated
$dq$ maps while retaining known magnetic structure.

We begin with terminal voltage, current, and electrical speed. The stationary
voltage balance yields flux samples. A conservative magnetic description
then interprets those vectors as local slopes of a latent scalar co-energy
surface. Reflection symmetry constrains that surface, its gradient supplies
both flux components, and its Hessian supplies their differential inductances.
Energy-based machine descriptions and saturation models support this approach
\cite{jebai2014,jebai2011}; standard voltage conventions follow
\cite{binder2017,schroeder2021}.

This paper owns the full co-energy, symmetry, integrability, and fitting
development. Its companion studies the complementary question: how do
departures from this admissible structure identify resistance, rotor-angle,
and related reconstruction errors \cite{companion-diagnostics}? Here those
errors motivate preserving fit residuals; they are not estimated. Both papers
remain independently readable and carry their own framework and numerical
data. The synthetic evaluation below concerns noise, sampling, and derivative
quality rather than parameter-error injections.

Unadorned quantities denote a true constitutive field, measured electrical
inputs are explicitly marked, $\flux^{\rm rec}$ denotes reconstructed samples,
and $\hat\flux_{\rm fit}=\nabla\hat\coenergy$ denotes a fitted field. All magnetic
statements refer to fixed temperature and magnet state, with an effective
fundamental or consistently averaged quasi-static description. No additional
energy measurement is assumed.
\begin{figure}[tb]
\centering
\begin{tikzpicture}[node distance=8mm,]
 \node[pipeline block] (m) {Measured $u_d,u_q,i_d,i_q,\omega_e$\\and used $R_s$};
 \node[pipeline block,below=of m] (f) {Stationary reconstruction\\$\vect\psi^{\rm rec}(i_d,i_q)$ samples};
 \node[pipeline block,below=of f] (energyfit) {Symmetry and integrability; gradient fit\\latent $W'(i_d,i_q)$};
 \node[pipeline block,below=of energyfit] (d) {Analytic derivatives\\$\hat\flux_{\rm fit}=\nabla\hat W'$, $\hat\Ldiff=\nabla^2\hat W'$};
 \node[diagnostic block,text width=39mm,right=12mm of energyfit] (r)
  {Structured fit residual\\$\flux^{\rm rec}-\hat\flux_{\rm fit}$\\[2pt]
   Companion: error hypotheses\\and correction proposals};
 \draw[direction arrow] (m)--(f);\draw[direction arrow] (f)--(energyfit);
 \draw[direction arrow] (energyfit)--(d);\draw[direction arrow] (energyfit)--(r);
\end{tikzpicture}
\caption{Measurement-to-model pipeline. No energy sensor is required:
co-energy is inferred from the flux vector field. The diagnostic branch
preserves the unexplained residual instead of silently absorbing it.}
\label{fig:pipeline}
\end{figure}

